\documentclass[a4paper,11pt]{article}
        \usepackage[top=15mm,bottom=18mm,left=16mm,right=16mm]{geometry}
        \usepackage{fontspec}
        \usepackage{xeCJK}
        \setmainfont{Times New Roman}
        \setCJKmainfont{Yu Gothic}
        \setmonofont{Consolas}
        \setCJKmonofont{Yu Gothic}
        \usepackage{booktabs}
        \usepackage{longtable}
        \usepackage{tabularx}
        \usepackage{array}
        \usepackage{enumitem}
        \usepackage{hyperref}
        \usepackage{listings}
        \usepackage{xcolor}
        \hypersetup{
          colorlinks=true,
          linkcolor=blue,
          urlcolor=blue,
          pdftitle={live forecast 取得ノード},
          pdfauthor={Codex}
        }
        \lstset{
          basicstyle=\ttfamily\small,
          breaklines=true,
          columns=fullflexible,
          keepspaces=true,
          frame=single,
          framerule=0.2pt,
          rulecolor=\color{black},
          showstringspaces=false
        }
        \setlength{\parskip}{0.42em}
        \setlength{\parindent}{1em}
        \renewcommand{\arraystretch}{1.15}
        \title{22. live forecast 取得ノード}
        \author{solar\_ws0129-main}
        \date{2026-07-29}
        \begin{document}
        \maketitle

        \section*{対象}
        \begin{tabularx}{\linewidth}{>{\raggedright\arraybackslash}p{0.18\linewidth} X}
        \toprule
        項目 & 内容 \\
        \midrule
        ファイル & \path|mpc_solarcar/weather_fetch_node.py| \\
        種別 & Python \\
        区分 & runtime node \\
        役割 & chase GPS または fallback 座標から Open-Meteo forecast を取得し、planner が読む CSV を更新する。 \\
        起動文脈 & live 系の forecast 更新入口。 \\
        \bottomrule
        \end{tabularx}

        \section*{このファイルを読む理由}
        chase GPS または fallback 座標から Open-Meteo forecast を取得し、planner が読む CSV を更新する。

        \begin{itemize}[leftmargin=1.6em]
        \item raw forecast CSV を定期更新する。
\item planner 自体は topic ではなく forecast CSV を読む。
        \end{itemize}

        \section*{呼び出し関係}
        \subsection*{どこから呼ばれるか}
        \begin{itemize}[leftmargin=1.6em]
        \item \path|mpc_solarcar/live_launch.py|
        \end{itemize}

        \subsection*{次に読むと理解が進むファイル}
        \begin{itemize}[leftmargin=1.6em]
        \item \path|mpc_solarcar/weather_utils.py|
\item \path|mpc_solarcar/wind_correction_node.py|
        \end{itemize}

        \subsection*{同一リポジトリ内の主要依存}
        \begin{itemize}[leftmargin=1.6em]
        \item \path|mpc_solarcar/solar_profile.py|
\item \path|mpc_solarcar/weather_utils.py|
        \end{itemize}

        \section*{ソース冒頭の把握}
        モジュール docstring または先頭意図の要約:

        モジュール docstring は明示されていない。

        \subsection*{代表コード断片}
        \begin{lstlisting}[language=Python]
from .weather_utils import fetch_openmeteo_forecast, write_forecast_csv


class WeatherFetchNode(Node):
    def __init__(self) -> None:
        super().__init__("weather_fetch_node")
        self.declare_parameter("profile_yaml", "")
        self.declare_parameter("output_csv", "")
        self.declare_parameter("raw_forecast_csv", "")
        self.declare_parameter("gps_topic", "/chase/gps")
        self.declare_parameter("fetch_period_sec", 3600.0)
        self.declare_parameter("forecast_days", 3)
        self.declare_parameter("step_minutes", 10)
        self.declare_parameter("timezone_name", "Australia/Darwin")
        self.declare_parameter("fallback_latitude", -12.4634)
        self.declare_parameter("fallback_longitude", 130.8456)
        self.declare_parameter("tcell_gain", 0.03)

        profile_yaml = str(self.get_parameter("profile_yaml").value or "").strip()
        if profile_yaml:
            profile_path, cfg = load_profile(profile_yaml)
            runtime_cfg = get_section(cfg, "runtime")
\end{lstlisting}


        \section*{主要構造の読み取り}
        主要クラスは WeatherFetchNode。 主要関数は main。 ROS パラメータ宣言は 11 件。 ROS I/O は publisher 1 件、subscription 1 件。

        \subsection*{主要クラス・関数}
            \begin{longtable}{p{0.07\linewidth} p{0.16\linewidth} p{0.25\linewidth} p{0.40\linewidth}}
            \toprule
            行 & 種別 & 名前 & 補足 \\
            \midrule
            15 & class & \path|WeatherFetchNode| & bases: Node \\
16 & function & \path|__init__| & args: self \\
104 & function & \path|_on_gps| & args: self, msg \\
111 & function & \path|_fetch_once| & args: self \\
138 & function & \path|main| &  \\
            \bottomrule
            \end{longtable}

        \subsection*{ファイルを上から読んだときの定義順}
            \begin{longtable}{p{0.07\linewidth} p{0.84\linewidth}}
            \toprule
            行 & 内容 \\
            \midrule
            15 & クラス WeatherFetchNode を定義する。 \\
138 & 関数 main を定義する。 \\
            \bottomrule
            \end{longtable}

        \subsection*{import 群の詳細}
            \begin{longtable}{p{0.07\linewidth} p{0.33\linewidth} p{0.50\linewidth}}
            \toprule
            行 & import 文 & このファイルで読む理由 \\
            \midrule
            1 & \path|from __future__ import annotations| & 型ヒントの遅延評価を有効にし、自己参照型や forward reference を安全に扱うため。 このファイル内での使用位置は少ないか、間接利用である。 \\
3 & \path|from pathlib import Path| & ファイルやディレクトリを安全に扱うため。 このファイル内での主な使用位置は L85, L86。 \\
5 & \path|import rclpy| & ROS 2 Python ノードとして起動・spin するため。 このファイル内での主な使用位置は L39, L40, L44, L48, L52, L56, L60, L69, ...。 \\
6 & \path|from rclpy.node import Node| & ROS 2 ノード本体の基底クラスとして使うため。 このファイル内での主な使用位置は L15。 \\
8 & \path|from sensor_msgs.msg import NavSatFix| & GPS 緯度経度高度を ROS topic でやり取りするため。 このファイル内での主な使用位置は L99, L104。 \\
9 & \path|from std_msgs.msg import String| & Float32 や String など軽量メッセージで planner / vehicle 値を publish するため。 このファイル内での主な使用位置は L98, L130, L134。 \\
11 & \path|from .solar_profile import get_path, get_section, load_profile| & profile YAML 読込と検証 から get\_path, get\_section, load\_profile を読み込み、このファイルの内部処理を分担させるため。 実体ファイルは mpc\_solarcar/solar\_profile.py。 このファイル内での主な使用位置は L32, L33, L34, L35, L39。 \\
12 & \path|from .weather_utils import fetch_openmeteo_forecast, write_forecast_csv| & weather\_utils.py から fetch\_openmeteo\_forecast, write\_forecast\_csv を読み込み、このファイルの内部処理を分担させるため。 実体ファイルは mpc\_solarcar/weather\_utils.py。 このファイル内での主な使用位置は L115, L123, L128。 \\
            \bottomrule
            \end{longtable}

        \section*{関数・クラスを上から順に解説}
        \subsection*{L15 クラス \path|WeatherFetchNode|}
            \begin{itemize}[leftmargin=1.6em]
              \item 定義: \path|WeatherFetchNode(bases=Node)|
              \item docstring: 明示されていない。
              \item このブロックが直接呼ぶ主な関数・メソッド: Parameter, Path, String, \_\_init\_\_, \_fetch\_once, create\_publisher, create\_subscription, create\_timer, declare\_parameter, error, expanduser, fetch\_openmeteo\_forecast
              \item 戻り値の要点: 戻り値が無いか、return 文が明示されていない。
            \end{itemize}
            \begin{enumerate}[leftmargin=1.8em]
            \item 関数 \_\_init\_\_ を定義する。
\item 関数 \_on\_gps を定義する。
\item 関数 \_fetch\_once を定義する。
            \end{enumerate}

\subsection*{L138 関数 \path|main|}
            \begin{itemize}[leftmargin=1.6em]
              \item 定義: \path|main()|
              \item docstring: 明示されていない。
              \item このブロックが直接呼ぶ主な関数・メソッド: WeatherFetchNode, destroy\_node, init, shutdown, spin
              \item 戻り値の要点: 戻り値が無いか、return 文が明示されていない。
            \end{itemize}
            \begin{enumerate}[leftmargin=1.8em]
            \item rclpy.init(...) を実行する。
\item node に WeatherFetchNode() の結果を代入する。
\item rclpy.spin(...) を実行する。
\item node.destroy\_node(...) を実行する。
\item rclpy.shutdown(...) を実行する。
            \end{enumerate}

        \subsection*{設定値と入力パラメータ}
            \begin{longtable}{p{0.07\linewidth} p{0.35\linewidth} p{0.45\linewidth}}
            \toprule
            行 & 名前 & 既定値・補足 \\
            \midrule
            18 & \path|profile_yaml| &  \\
19 & \path|output_csv| &  \\
20 & \path|raw_forecast_csv| &  \\
21 & \path|gps_topic| & /chase/gps \\
22 & \path|fetch_period_sec| & 3600.0 \\
23 & \path|forecast_days| & 3 \\
24 & \path|step_minutes| & 10 \\
25 & \path|timezone_name| & Australia/Darwin \\
26 & \path|fallback_latitude| & -12.4634 \\
27 & \path|fallback_longitude| & 130.8456 \\
28 & \path|tcell_gain| & 0.03 \\
            \bottomrule
            \end{longtable}

        \subsection*{ROS topic I/O}
            \begin{longtable}{p{0.07\linewidth} p{0.18\linewidth} p{0.34\linewidth} p{0.28\linewidth}}
            \toprule
            行 & 種別 & topic & callback / 補足 \\
            \midrule
            98 & publisher & \path|/weather/fetch_status| & publisher \\
99 & subscription & \path|self.gps_topic| & self.\_on\_gps \\
            \bottomrule
            \end{longtable}







        \section*{処理の流れ}
        \begin{enumerate}[leftmargin=1.7em]
        \item 初期化時に設定値や入力パスを読み込む。
\item publisher / subscription / timer を準備する。
\item timer callback 周期で主処理を進める。
        \end{enumerate}

        \section*{このファイルの位置づけ}
        この資料群では、\path|mpc_solarcar/weather_fetch_node.py| を
        \textbf{runtime node}の中核として扱う。
        したがって、ここで扱う処理は単独で完結せず、
        前段の profile / launch / telemetry と後段の planner / logger / report へ繋がる。

        \end{document}
